\documentclass[a4paper]{ltjsreport}
\usepackage{luatexja-fontspec}
\usepackage{graphicx}
\usepackage{titlesec}
\usepackage{url}
\usepackage{comment}
\usepackage[hidelinks]{hyperref}
\setmainfont{Hiragino Mincho ProN}
\setsansfont{Hiragino Sans}
\setcounter{secnumdepth}{3} % subsubにも番号が振られる　目次には出ない

% 章タイトルの設定
\titleformat{\chapter}[display]
  {\normalfont\huge}
  {\chaptertitlename\ \thechapter 章}
  {20pt}
  {\LARGE}
% 節の設定 
\titleformat{\section}
  {\normalfont\Large}
  {\thesection}
  {1em}
  {}
% 項の設定 
\titleformat{\subsection}
  {\normalfont\large}
  {\thesubsection}
  {1em}
  {}
% 目次に出ない小見出しの設定
\titleformat{\subsubsection}
  {\normalfont\large}
  {\thesubsubsection}
  {1em}
  {}

\title{速報実証衛星ARICA-2における〇〇の開発}
\author{岩永　知沙季}
\date{\today}

\begin{document}

\maketitle

\tableofcontents % 目次を自動生成

\chapter{はじめに}

\section{研究背景}
\subsection{ガンマ線バースト（Gamma-Ray Burst：GRB）}
\subsection{ガンマ線バースト速報システムの現状}
\subsection{速報実証衛星ARICA（AGU Remote Innovative CubeSat Alert system）}

\section{本研究の目的}

\section{本論文の構成}
\chapter{速報実証衛星ARICA-2（AGU Remote Innovative CubeSat Alert system -2）}
\section{ARICA-2プロジェクトの概要}
\section{サクセスクライテリア}
\section{ARICA-2の構成}
\chapter{姿勢決定系・姿勢制御系}
\section{姿勢決定系}
\subsection{}
\subsection{姿勢決定系のコンポーネント}

\section{姿勢制御系}
\subsection{}
\subsection{姿勢制御系のコンポーネント}


\chapter{基板の開発}
\section{概要}
衛星システムは複数の電子回路基板によって構成され、それらが電気的に結合することで機能する。本開発は、機能実証を行うための試作機であるエンジニアリングモデル（EM）の開発と、実機搭載用のフライトモデル（FM）の開発の大きく分けて二段階を経て実施された。以下に、開発した基板とその開発プロセスの概要を述べる。

\subsection{基板の概要}
磁気トルカボード、（コネクタボード）、姿勢系システム搭載ボード、赤外線センサインターフェースボード、Spresense拡張ボード

\subsection{基板設計ソフトウェア EAGLE}
本研究の基板開発では、Autodesk社が提供するプリント基板設計者向けの電子設計オートメーションソフトウェア「EAGLE」を用いた。EAGLEで基板設計を行った後、そのガーバーデータを基板製造会社であるP板.comに提出して基板製造を依頼することで、自身の設計した基板が完成する。

Autodesk社のEAGLEは2026年6月7日をもって販売およびサポートを終了することが決定しており、その後は同社Autodesk Fusion内のAutodesk Fusion Electronicsで同様の基板設計ができる。本研究で開発した主要な基板のEAGLEデータについては、Autodesk Fusion Electronicsに移行して〇〇に置いてある。% memo 保存場所を書く

\subsection{開発プロセスの概要}
FM基板の開発には、EM基板の課題を受けた設計変更と要求
基板の開発は、設計要求の確認、基板仕様の検討、回路設計、アートワーク設計、基板の評価という順で行う。以下にそれぞれの詳細をまとめる。

\subsubsection{設計要求の確認}
\subsubsection{基板仕様の検討}
\subsubsection{回路設計}
\subsubsection{アートワーク設計}
\subsubsection{基板の評価}

\section{エンジニアリングモデルの基板開発}
\subsection{基板仕様}
\subsection{課題とその究明}


\section{フライトモデルの基板開発}
\subsection{EMの課題を受けて}
\subsection{設計要求}
\subsubsection{磁気トルカボード}
\subsubsection{磁気トルカコネクタボード}
\subsubsection{姿勢系システム搭載ボード}
\subsubsection{赤外線センサインターフェースボード}
\subsubsection{Spresense拡張ボード}


\subsection{基板仕様の検討}
\subsubsection{磁気トルカボード}
\subsubsection{磁気トルカコネクタボード}
\subsubsection{姿勢系システム搭載ボード}
\subsubsection{赤外線センサインターフェースボード}
\subsubsection{Spresense拡張ボード}

\subsection{回路設計}
\subsubsection{磁気トルカボード}
\subsubsection{磁気トルカコネクタボード}
\subsubsection{姿勢系システム搭載ボード}
\subsubsection{赤外線センサインターフェースボード}
\subsubsection{Spresense拡張ボード}

\subsection{アートワーク設計}
\subsubsection{磁気トルカボード}
\subsubsection{磁気トルカコネクタボード}
\subsubsection{姿勢系システム搭載ボード}
\subsubsection{赤外線センサインターフェースボード}
\subsubsection{Spresense拡張ボード}

\subsection{基板の評価}
\subsubsection{磁気トルカボード}
\subsubsection{磁気トルカコネクタボード}
\subsubsection{姿勢系システム搭載ボード}
\subsubsection{赤外線センサインターフェースボード}
\subsubsection{Spresense拡張ボード}

\section{（仮称）ハードウェア設計の注意点}



% memo ハーネスのことも述べる
\chapter{赤外線センサ関連}

\chapter{おわりに}

\section{まとめ}
\section{今後の展望}

\include{chapters/thanks}

\appendix
\chapter{基板設計の手引き}

\chapter{シリアル通信とハードウェア設計の}

\begin{thebibliography}{99}
\bibitem{1} 河合誠之, 浅野勝晃, ガンマ線バースト, 日本評論社, 2019, 新天文学ライブラリー第5巻
\bibitem{2} Jeanette Kazmierczak. “NASA’s Swift, Fermi Missions Detect Exceptional Cosmic Blast”. NASA(.gov). 2023-07-26. \url{https://www.nasa.gov/universe/nasas-swift-fermi-missions-detect-exceptional-cosmic-blast/}, (参照日：2024-2-20)
\bibitem{3} 坂本貴紀, HETE-2って何？, \url{http://www.hp.phys.titech.ac.jp/hete2/whatishete.html}, (参照日：2024-2-20)
\bibitem{4} Lorella Angelini Jesse Allen, The HETE-2 Satellite, 2020-09-24, \url{https://heasarc.gsfc.nasa.gov/docs/hete2/hete2.html}, (2024-02-17)
\bibitem{5} 児玉太志, 民間通信衛星を用いたガンマ線バースト速報システムの開発, 2018年度 修士論文
\bibitem{6} 一般財団法人リモート・センシング技術センター, TDRS, \url{https://www.restec.or.jp/satellite/tdrs.html}
\bibitem{7} J.D. Myers, The Neil Gehrels Swift Observatory, 2023-09-07, \url{https://swift.gsfc.nasa.gov/}, (2024-02-18)
\bibitem{8} KDDI株式会社, 【KDDI】イリジウムサービス｜衛星通信/BCP対策｜法人向け, \url{https://biz.kddi.com/service/iridium/}, (参照日：2024-2-20)
\bibitem{9} Julien's Lab, About｜Iridium Satellites Live Map, \url{https://iridiumwhere.com/about/}, (参照日：2024-2-20)
\bibitem{10} Globalstar, Satellite Solutions and Services｜Globalstar US, \url{https://www.globalstar.com/en-us/}, (参照日：2024-2-20)
\bibitem{11} ソニーセミコンダクタソリューションズ株式会社, IoT用ボードコンピュータSPRESENSE™, \url{https://www.sony-semicon.com/ja/products/spresense/index.html}, (参照日：2024-2-20)
\bibitem{12} JAXA, 革新的衛星技術実証２号機, \url{https://www.kenkai.jaxa.jp/kakushin/kakushin02.html}, (参照日：2024-2-20)
\bibitem{13} Omron Corporation, 形 D6T MEMS 非接触温度センサ, \url{https://components.omron.com/jp-ja/datasheet_pdf/CDSC-012.pdf}, (参照日：2024-2-20)
\bibitem{14} Pimoroni, ICM20948 9DoF Motion Sensor Breakout, \url{https://shop.pimoroni.com/products/icm20948?variant=27843993960531}, (参照日：2024-2-20)
\bibitem{15} 株式会社秋月電子通商, DRV8830モータードライバモジュール, \url{https://akizukidenshi.com/catalog/g/g106273/}, (参照日：2024-2-20)
\bibitem{16} 辻祐樹, 速報実証衛星 ARICA-2 において用いる太陽センサの開発及び I2C 通信を用いるデバイスの統合, 2022年度 学士論文
\bibitem{17} 東阪電子機器株式会社, I2C(Inter-Integrated Circuit) について, \url{https://www.lapis-tech.com/jp/common/miconlp/tips/2-17-i2cとは/}, (参照日：2024-2-20)
\bibitem{18} 青山学院大学理工学部物理・数理学科 坂本研究室, ARICA-2 Preliminary Design Review, \url{https://docs.google.com/presentation/d/16HseKSFk4e2CegcKgMQAkCcH3iP3NXOl/edit#slide=id.p131}, (参照日：2024-2-20)
\bibitem{19} CubeSpace Satellite System, CubeTorquer Generation 2 Interface Control Document(ICD), \url{https://www.cubespace.co.za/downloads/cs-dev.icd.cr-01_cubetorquer_gen2_icd_ver.1.00.pdf}, (参照日：2024-2-20)
\bibitem{20} 湯本電機株式会社, PEEKとは？高機能の熱可塑性樹脂の特徴と各種グレード, \url{https://www.yumoto.jp/material-onepoint/plastic-peek}, (参照日：2024-2-20)
\end{thebibliography}


\begin{comment}

・「memo」で検索すると修正箇所がある
・改行は\\ ただし字下げはされない
・新しい段落にする際は一行空白を入れる　字下げもされる
・*をつけると目次に表示されない

・画像テンプレ

\begin{figure}[htbp]
  \centering
  \includegraphics[width=0.8\textwidth]{images/スクリーンショット 2025-08-03 17.23.28.png}
  \caption{ここに図のキャプション（説明文）を書きます。}
  \label{fig:myfigure}
\end{figure}

図\ref{fig:myfigure}は、本研究の概念図です。

\end{comment}

\end{document}